% =============== Introduction =================================================
% Goal: 1.5-2 pages · structure:
%   (1) the problem · (2) why existing solutions don't fit
%   (3) our approach · (4) contributions
% =============================================================================

Modern coding agents based on Large Language Models (LLMs) such as
Claude Code, Cursor, and Continue.dev sustain dialogue sessions
spanning hundreds to thousands of turns. Within such sessions, a
recurring failure mode is \emph{persona drift}: the assistant
gradually forgets user-specified constraints (``always verify before
claiming success''), slips into bad habits the user explicitly
flagged earlier (``don't fabricate prior agreements''), or confabulates
having reached agreements that were never made. These behaviors are
not failures of factual recall---retrieval-augmented memory plugins
\citep{mem0,letta} successfully retrieve the relevant memos---but
failures of \emph{behavioral consistency}: the assistant has the
information but does not act on it.

\paragraph{Existing approaches fall into two categories.}
\textit{(a) Memory plugins} like mem0 \citep{mem0}, Letta
\citep{letta}, and Zep \citep{zep} focus on retrieval quality:
finding the right memory entries to inject into the prompt. While
these systems produce strong retrieval metrics (e.g., on
LongMemEval \citep{longmemeval2024}), they leave open the question
of whether the LLM \emph{honors} the retrieved content.
\textit{(b) White-box safety methods}, exemplified by
Persona Vectors \citep{personavectors2025}, identify directions in
the model's activation space corresponding to specific traits and
enable monitoring or steering of trait shifts. However, these
methods require access to model weights or hidden states, putting
them out of reach of the typical end-user who interacts with
proprietary LLMs through black-box APIs.

\paragraph{Our contribution.} We present \textbf{zenmind-mem}, a
black-box persona drift detection system that runs entirely in user-space
hooks of an LLM coding agent (we target Claude Code in our reference
implementation, but the design is portable). At each user-issued
prompt, we compute the cosine similarity between the prompt and a
small set of \emph{behavioral anchor texts}, aggregated via weighted
top-$k$ mean. Anchors come in two flavors: positive (desired task
patterns) and negative (mistakes the user has flagged). The
difference between aggregated positive and negative similarity
yields a continuous \emph{drift score}, which we threshold into a
three-band output (\emph{aligned} / \emph{neutral} /
\emph{deviation}) and inject into the LLM's context for the current
turn.

\paragraph{Methodological insight.}
Our central finding is that the design of the anchor set, not the
choice of embedder or aggregation function, is the dominant factor
determining detection accuracy. We document a four-step iteration
that improves ROC AUC from 0.51 (random) to 0.92 on a 100-prompt
synthetic test set, and report negative results (e.g., that
substituting a cross-encoder reranker for the bi-encoder cosine
yields no measurable improvement at $36\times$ the latency).

\paragraph{Open source and reproducibility.}
All code, anchor sets, evaluation scripts, and benchmark data are
released under MIT (anchors under CC0) at
\url{https://github.com/chunxiaoxx/zenmind-mem}. The four-step
ablation, the head-to-head comparison with mem0 on LongMemEval-S,
and the negative cross-encoder result are all reproducible from
the repository in under one hour on a CPU-only machine.

\paragraph{Relation to Persona Vectors.}
Our work is \emph{complementary, not competing}, with
\citet{personavectors2025}. Persona Vectors operates in
activation space (white-box), giving more fine-grained signal but
requiring model weights. zenmind-mem operates in prompt-text space
(black-box), giving coarser signal but deployable by anyone with
hook-level access to the agent. We argue that a robust safety
posture for production LLM agents will combine both layers.
